%=============================================================
\chapter{Présentation des résultats}
%=============================================================

Ce chapitre présente l'ensemble des résultats obtenus dans le cadre
de ce travail, structurés autour de deux axes complémentaires. Le
premier axe est consacré à la présentation et à l'interprétation des
résultats issus de l'analyse statistique des données de l'enquête
réalisée par le \gls{minepat} sur la cartographie et les besoins des
entreprises du tissu productif local. Cette analyse s'articule en
trois niveaux de profondeur croissante~: une analyse univariée
décrivant la distribution des entreprises selon leurs caractéristiques
principales, une analyse bivariée croisant ces caractéristiques avec
les besoins exprimés et validant les associations observées par le
test du $\chi^2$ de Pearson, et une analyse multidimensionnelle
mobilisant l'Analyse des Correspondances Multiples (ACM), la
Classification Ascendante Hiérarchique (CAH) et l'Analyse Factorielle
Discriminante (AFD). Le second axe est consacré à la présentation de
la plateforme \textit{Business Partnership}, développée pour répondre
aux besoins identifiés lors de cette analyse. L'ensemble de ces
analyses ont été conduites sur un fichier contenant
\textbf{1\,194 entreprises}, issues de l'enquête nationale du
\gls{minepat} de 2025.

%=============================================================
\section{Résultats sur la cartographie des entreprises}
\label{sec:resultats_carto}
%=============================================================

Cette section présente les résultats de l'analyse statistique des
données collectées lors de l'enquête. Elle est structurée selon une
progression méthodologique allant du simple au complexe~: description
univariée, analyse des associations bivariées, puis exploration
multidimensionnelle.

%-------------------------------------------------------------
\subsection{Résultats de l'analyse univariée}
\label{subsec:univariee}
%-------------------------------------------------------------

L'analyse univariée a consisté à déterminer le nombre et le
pourcentage d'entreprises présentant des modalités similaires sur
les différentes variables retenues~: statut juridique, région
d'implantation, taille, ancienneté et secteur d'activité principal.
Elle constitue la première photographie descriptive du tissu productif
camerounais couvert par l'enquête, et permet d'identifier les
caractéristiques dominantes des entreprises interrogées avant
d'explorer les relations entre leurs attributs et leurs besoins.

\subsubsection{Répartition des entreprises par statut juridique}

La distribution des entreprises selon leur statut juridique révèle
une prépondérance des Sociétés à Responsabilité Limitée (SARL), qui
représentent 35,0\,\% des entreprises de l'échantillon. Ce résultat
reflète la tendance observée dans l'ensemble du tissu économique
camerounais, où la SARL constitue la forme juridique la plus courante
pour les entreprises de taille intermédiaire, en raison notamment de
la simplicité relative de sa constitution et de la limitation de
responsabilité qu'elle offre à ses associés.

\begin{figure}[H]
	\centering
	\includegraphics[width=0.75\textwidth]{./img/repartition_statut_juridique.png}
	\caption{Répartition des entreprises enquêtées par statut juridique}
	\label{fig:statut_juridique}
\end{figure}

\subsubsection{Répartition des entreprises par région}

\begin{figure}[H]
	\centering
	\includegraphics[width=0.75\textwidth]{./img/repartition_region.png}
	\caption{Répartition des entreprises enquêtées par région}
	\label{fig:region}
\end{figure}

La répartition géographique des entreprises enquêtées met en évidence
une distribution couvrant l'ensemble des dix régions du Cameroun.
[Compléter avec les chiffres issus de la figure~\ref{fig:region}
et commenter les régions les plus représentées.]

\subsubsection{Répartition des entreprises par taille}

\begin{figure}[H]
	\centering
	\includegraphics[width=0.75\textwidth]{./img/repartition_taille.png}
	\caption{Répartition des entreprises enquêtées par taille}
	\label{fig:taille}
\end{figure}

La distribution par taille reflète la structure du tissu productif
camerounais, caractérisé par une forte prédominance des très petites
et petites entreprises, conformément aux résultats du Recensement
Général des Entreprises de 2023 qui indiquait que 92,4\,\% des unités
économiques sont des entreprises individuelles \citep{ins_rge_2023}.
[Compléter avec les chiffres issus de la figure~\ref{fig:taille}.]

\subsubsection{Répartition des entreprises par ancienneté}

\begin{figure}[H]
	\centering
	\includegraphics[width=0.75\textwidth]{./img/repartition_anciennete.png}
	\caption{Répartition des entreprises enquêtées par ancienneté}
	\label{fig:anciennete}
\end{figure}

[Commenter la distribution par ancienneté en lien avec la figure
\ref{fig:anciennete}, en soulignant notamment la part des entreprises
jeunes (moins de 5 ans) et celle des entreprises matures (plus de
10 ans), et en interprétant ce que cela dit sur la dynamique
entrepreneuriale nationale.]

\subsubsection{Répartition des entreprises par secteur d'activité}

\begin{figure}[H]
	\centering
	\includegraphics[width=0.75\textwidth]{./img/repartition_secteur.png}
	\caption{Répartition des entreprises enquêtées par secteur d'activité}
	\label{fig:secteur}
\end{figure}

[Commenter la distribution sectorielle. Souligner la prédominance
attendue du secteur tertiaire, conformément aux résultats du RGE3 de
l'INS qui indiquait que 85,6\,\% des unités économiques relèvent du
secteur tertiaire \citep{ins_rge_2023}.]

\subsubsection{Répartition des besoins exprimés}

\begin{figure}[H]
	\centering
	\includegraphics[width=0.75\textwidth]{./img/repartition_besoins.png}
	\caption{Répartition globale des besoins exprimés par les entreprises
		enquêtées}
	\label{fig:besoins_global}
	\end{figure>
		
		L'analyse de la distribution globale des besoins exprimés par les
		entreprises met en évidence que le besoin en équipements et le besoin
		en financement constituent les deux besoins les plus fréquemment
		cités par les répondants. Ce résultat est cohérent avec les constats
		établis lors des enquêtes antérieures du \gls{minepat}, qui avaient
		mis en évidence que l'accès aux facteurs de production et aux
		financements constitue le principal frein au développement des
		entreprises camerounaises \citep{minepat_enquete_2025}. Le besoin en
		innovation et en recherche, bien que moins fréquemment cité, demeure
		non négligeable et témoigne d'une prise de conscience progressive des
		entreprises quant à l'importance de la modernisation technologique
		comme levier de compétitivité.
		
		%-------------------------------------------------------------
		\subsection{Résultats de l'analyse bivariée}
		\label{subsec:bivariee}
		%-------------------------------------------------------------
		
		L'analyse bivariée a consisté à croiser les variables descriptives
		des entreprises avec les besoins qu'elles ont exprimés, afin
		d'identifier d'éventuelles associations structurelles entre les
		caractéristiques des entreprises et leurs besoins spécifiques. Deux
		outils ont été mobilisés~: les tableaux de contingence, qui permettent
		de visualiser la distribution conjointe des variables, et le test du
		$\chi^2$ de Pearson accompagné du coefficient $V$ de Cramér, qui
		permettent de tester et de quantifier l'intensité de ces associations
		\citep{saporta_2010}.
		
		\subsubsection{Répartition des besoins selon le statut juridique}
		
		Les figures ci-dessous présentent la distribution des besoins des
		entreprises selon leur statut juridique, ainsi que les résultats du
		test d'association correspondant.
		
		\begin{figure}[H]
			\centering
			\includegraphics[width=\textwidth]{./img/besoins_statut.png}
			\caption{Distribution des besoins des entreprises selon le statut
				juridique}
			\label{fig:besoins_statut}
		\end{figure}
		
		\begin{table}[H]
			\centering
			\caption{Test du $\chi^2$ entre le statut juridique et les besoins
				des entreprises}
			\label{tab:khi2_statut}
			\renewcommand{\arraystretch}{1.4}
			\begin{tabularx}{\textwidth}{>{\raggedright\arraybackslash}X
					r r
					>{\centering\arraybackslash}p{2.8cm}
					r}
				\toprule
				\rowcolor{primary!15}
				\textbf{Besoin} & \textbf{$\chi^2$} & \textbf{p-valeur}
				& \textbf{Significatif} & \textbf{$V$ Cramér} \\
				\midrule
				Besoin en intrants        & 14,26 & 0,027  & Oui ($p < 0{,}05$) & 0,109 \\
				\rowcolor{lightgray}
				Besoin en financement     &  2,94 & 0,817  & Non                & 0,050 \\
				Besoin en emballages      & 44,47 & $<0{,}001$ & Oui ($p < 0{,}05$) & 0,193 \\
				\rowcolor{lightgray}
				Besoin en transport       & 50,44 & $<0{,}001$ & Oui ($p < 0{,}05$) & 0,206 \\
				Appui aux entreprises     & 26,32 & $<0{,}001$ & Oui ($p < 0{,}05$) & 0,148 \\
				\rowcolor{lightgray}
				Besoin en équipements     & 61,66 & $<0{,}001$ & Oui ($p < 0{,}05$) & 0,227 \\
				Innovation et recherche   & 77,88 & $<0{,}001$ & Oui ($p < 0{,}05$) & 0,255 \\
				\bottomrule
			\end{tabularx}
		\end{table}
		
		Les résultats du test du $\chi^2$ révèlent que le statut juridique
		est significativement associé à six des sept besoins étudiés, à
		l'exception du besoin en financement ($\chi^2 = 2{,}94$~;
		$p = 0{,}817$). Les associations les plus fortes concernent le besoin
		en innovation et en recherche ($V = 0{,}255$) et le besoin en
		équipements ($V = 0{,}227$), suggérant que la nature juridique de
		l'entreprise influence significativement la manière dont elle exprime
		ses besoins technologiques. Les entreprises de type SUARL/SARLU
		présentent notamment le taux le plus élevé de besoin en appui aux
		entreprises (20,0\,\%) et en financement (26,2\,\%), tandis que les
		Coopératives expriment davantage de besoins en innovation
		(14,1\,\%). Le besoin en financement constitue une exception notable~:
		il est exprimé de manière relativement homogène quelle que soit la
		forme juridique de l'entreprise, ce qui suggère que la contrainte
		financière est une réalité transversale à l'ensemble du tissu
		économique, indépendante du cadre juridique de l'entreprise.
		
		\subsubsection{Répartition des besoins selon la région}
		
		\begin{figure}[H]
			\centering
			\includegraphics[width=\textwidth]{./img/besoins_region.png}
			\caption{Distribution des besoins des entreprises selon la région
				d'implantation}
			\label{fig:besoins_region}
		\end{figure}
		
		\begin{table}[H]
			\centering
			\caption{Test du $\chi^2$ entre la région d'implantation et les besoins
				des entreprises}
			\label{tab:khi2_region}
			\renewcommand{\arraystretch}{1.4}
			\begin{tabularx}{\textwidth}{>{\raggedright\arraybackslash}X
					r r
					>{\centering\arraybackslash}p{2.8cm}
					r}
				\toprule
				\rowcolor{primary!15}
				\textbf{Besoin} & \textbf{$\chi^2$} & \textbf{p-valeur}
				& \textbf{Significatif} & \textbf{$V$ Cramér} \\
				\midrule
				Besoin en intrants        & 28,91 & $<0{,}001$ & Oui ($p < 0{,}05$) & 0,156 \\
				\rowcolor{lightgray}
				Besoin en financement     &  6,21 & 0,718  & Non                    & 0,072 \\
				Besoin en emballages      & 18,69 & 0,028  & Oui ($p < 0{,}05$)    & 0,125 \\
				\rowcolor{lightgray}
				Besoin en transport       & 29,83 & $<0{,}001$ & Oui ($p < 0{,}05$) & 0,158 \\
				Appui aux entreprises     & 23,36 & 0,005  & Oui ($p < 0{,}05$)    & 0,140 \\
				\rowcolor{lightgray}
				Besoin en équipements     & 28,25 & $<0{,}001$ & Oui ($p < 0{,}05$) & 0,154 \\
				Innovation et recherche   & 23,77 & 0,005  & Oui ($p < 0{,}05$)    & 0,141 \\
				\bottomrule
			\end{tabularx}
		\end{table}
		
		La région d'implantation s'avère significativement associée à six
		besoins sur sept, le besoin en financement étant à nouveau le seul
		pour lequel aucune association significative n'est détectée
		($\chi^2 = 6{,}21$~; $p = 0{,}718$). Ce résultat confirme le
		caractère transversal de la contrainte financière, commune à
		l'ensemble des régions. Les besoins en transport ($V = 0{,}158$)
		et en intrants ($V = 0{,}156$) présentent les associations les plus
		fortes avec la région, ce qui est cohérent avec les disparités
		infrastructurelles importantes entre les régions camerounaises.
		Les régions du Nord-Ouest et de l'Est, qui concentrent une part
		significative des entreprises à besoins en transport élevés,
		illustrent comment les contraintes géographiques et l'état des
		infrastructures conditionnent les priorités des entreprises locales.
		
		\subsubsection{Répartition des besoins selon l'ancienneté}
		
		\begin{figure}[H]
			\centering
			\includegraphics[width=\textwidth]{./img/besoins_anciennete.png}
			\caption{Distribution des besoins des entreprises selon leur
				ancienneté}
			\label{fig:besoins_anciennete}
		\end{figure}
		
		\begin{table}[H]
			\centering
			\caption{Test du $\chi^2$ entre l'ancienneté et les besoins
				des entreprises}
			\label{tab:khi2_anciennete}
			\renewcommand{\arraystretch}{1.4}
			\begin{tabularx}{\textwidth}{>{\raggedright\arraybackslash}X
					r r
					>{\centering\arraybackslash}p{2.8cm}
					r}
				\toprule
				\rowcolor{primary!15}
				\textbf{Besoin} & \textbf{$\chi^2$} & \textbf{p-valeur}
				& \textbf{Significatif} & \textbf{$V$ Cramér} \\
				\midrule
				Besoin en intrants        &  6,46 & 0,264  & Non                 & 0,074 \\
				\rowcolor{lightgray}
				Besoin en financement     &  4,44 & 0,488  & Non                 & 0,061 \\
				Besoin en emballages      &  5,09 & 0,405  & Non                 & 0,065 \\
				\rowcolor{lightgray}
				Besoin en transport       &  5,52 & 0,356  & Non                 & 0,068 \\
				Appui aux entreprises     & 11,84 & 0,037  & Oui ($p < 0{,}05$)  & 0,100 \\
				\rowcolor{lightgray}
				Besoin en équipements     & 14,84 & 0,011  & Oui ($p < 0{,}05$)  & 0,111 \\
				Innovation et recherche   & 20,10 & 0,001  & Oui ($p < 0{,}05$)  & 0,130 \\
				\bottomrule
			\end{tabularx}
		\end{table}
		
		L'ancienneté de l'entreprise est le critère présentant le moins
		d'associations significatives avec les besoins, puisque seulement
		trois besoins sur sept lui sont significativement liés. Les entreprises
		jeunes (moins de 3 ans) expriment davantage de besoins en financement
		(22,2\,\%) et en équipements (21,1\,\%), ce qui s'explique par leur
		stade de développement~: en phase de démarrage, l'accès aux
		ressources matérielles et financières constitue la contrainte
		prioritaire. À l'inverse, les entreprises plus matures (20 ans et
		plus) expriment un besoin en innovation et en recherche plus marqué
		(13,9\,\%), témoignant d'une préoccupation croissante pour la
		compétitivité à long terme et la modernisation technologique au fur
		et à mesure que l'entreprise se consolide.
		
		\subsubsection{Répartition des besoins selon le secteur d'activité}
		
		\begin{figure}[H]
			\centering
			\includegraphics[width=\textwidth]{./img/besoins_secteur.png}
			\caption{Distribution des besoins des entreprises selon le secteur
				d'activité}
			\label{fig:besoins_secteur}
		\end{figure}
		
		\begin{table}[H]
			\centering
			\caption{Test du $\chi^2$ entre le secteur d'activité et les besoins
				des entreprises}
			\label{tab:khi2_secteur}
			\renewcommand{\arraystretch}{1.4}
			\begin{tabularx}{\textwidth}{>{\raggedright\arraybackslash}X
					r r
					>{\centering\arraybackslash}p{2.8cm}
					r}
				\toprule
				\rowcolor{primary!15}
				\textbf{Besoin} & \textbf{$\chi^2$} & \textbf{p-valeur}
				& \textbf{Significatif} & \textbf{$V$ Cramér} \\
				\midrule
				Besoin en intrants        & 103,48 & $<0{,}001$ & Oui ($p < 0{,}05$) & 0,294 \\
				\rowcolor{lightgray}
				Besoin en financement     &  34,10 & 0,003  & Oui ($p < 0{,}05$)    & 0,169 \\
				Besoin en emballages      & 306,81 & $<0{,}001$ & Oui ($p < 0{,}05$) & 0,507 \\
				\rowcolor{lightgray}
				Besoin en transport       & 147,87 & $<0{,}001$ & Oui ($p < 0{,}05$) & 0,352 \\
				Appui aux entreprises     &  20,13 & 0,167  & Non                    & 0,130 \\
				\rowcolor{lightgray}
				Besoin en équipements     & 239,10 & $<0{,}001$ & Oui ($p < 0{,}05$) & 0,447 \\
				Innovation et recherche   & 163,49 & $<0{,}001$ & Oui ($p < 0{,}05$) & 0,370 \\
				\bottomrule
			\end{tabularx}
		\end{table}
		
		Le secteur d'activité constitue la variable présentant les
		associations les plus fortes avec les besoins des entreprises~:
		six besoins sur sept lui sont significativement liés, avec des
		valeurs de $V$ de Cramér particulièrement élevées, notamment pour
		le besoin en emballages ($V = 0{,}507$), le besoin en équipements
		($V = 0{,}447$) et le besoin en transport ($V = 0{,}352$). Ces
		résultats confirment l'intuition selon laquelle le secteur d'activité
		est le principal déterminant de la nature des besoins exprimés par
		les entreprises. Les entreprises du secteur des Services immobiliers
		expriment le taux de besoin en financement le plus élevé (50,0\,\%),
		tandis que les entreprises du secteur des Médias et Divertissement
		présentent un besoin en équipements particulièrement fort (40,0\,\%).
		Le besoin en appui aux entreprises est le seul besoin pour lequel
		aucune association significative avec le secteur d'activité n'est
		détectée ($p = 0{,}167$), confirmant son caractère transversal à
		l'ensemble des secteurs.
		
		\begin{remarque}
			Deux constats transversaux émergent de l'ensemble des analyses
			bivariées~: \textbf{(i)} le besoin en financement est le seul besoin
			systématiquement non associé aux variables descriptives (statut
			juridique, région, ancienneté), à l'exception du secteur d'activité,
			ce qui en fait une contrainte universelle du tissu économique
			camerounais, indépendante des caractéristiques structurelles des
			entreprises~; \textbf{(ii)} le secteur d'activité est la variable
			dont le pouvoir explicatif est le plus élevé, avec les valeurs de
			$V$ de Cramér les plus importantes, ce qui justifie son rôle central
			dans la définition de l'algorithme de mise en relation de la
			plateforme \textit{Business Partnership}.
		\end{remarque}
		
		%-------------------------------------------------------------
		\subsection{Résultats de l'analyse multidimensionnelle}
		\label{subsec:multi}
		%-------------------------------------------------------------
		
		Les analyses univariée et bivariée ont permis de décrire le tissu
		économique et d'identifier des associations entre les caractéristiques
		des entreprises et leurs besoins. Cependant, ces méthodes ne
		permettent pas d'appréhender simultanément l'ensemble des relations
		entre les variables dans leur globalité. L'analyse multidimensionnelle
		vient combler cette limite en offrant une représentation synthétique
		des données dans un espace factoriel de dimension réduite, à partir
		duquel des groupes d'entreprises aux profils similaires peuvent être
		construits et caractérisés. Elle s'articule en trois étapes
		successives~: l'ACM pour la réduction de dimension, la CAH pour
		la construction des classes, et l'AFD pour leur validation et
		caractérisation.
		
		\subsubsection{Analyse des Correspondances Multiples (ACM)}
		
		L'ACM a été conduite à l'aide du package \texttt{FactoMineR} du
		langage R \citep{lebart_morineau_piron_1995}, sur l'ensemble des
		variables qualitatives retenues pour décrire les entreprises et
		leurs besoins. Son objectif principal est d'identifier les axes
		factoriels sur lesquels les données présentent la plus grande
		variabilité, afin de produire une représentation synthétique du
		nuage d'entreprises dans un espace de dimension réduite.
		
		\paragraph{Valeurs propres et inertie expliquée}
		
		Le choix du nombre d'axes à retenir s'effectue à partir de
		l'examen de l'éboulis des valeurs propres, qui représente la
		part d'inertie totale expliquée par chacun des axes factoriels.
		
		\begin{figure}[H]
			\centering
			\includegraphics[width=0.75\textwidth]{./img/acm_valeurs_propres.png}
			\caption{Éboulis des valeurs propres de l'ACM --- part d'inertie
				expliquée par axe factoriel}
			\label{fig:acm_vp}
		\end{figure}
		
		[Commenter la figure~\ref{fig:acm_vp} en indiquant le pourcentage
		d'inertie expliqué par les deux ou trois premiers axes, et justifier
		le choix du nombre d'axes retenus pour la suite de l'analyse.
		Exemple type de commentaire~: \textit{Les deux premiers axes
			factoriels expliquent conjointement X\,\% de l'inertie totale, ce
			qui est suffisant pour capturer la structure principale des données
			et justifie leur retenue pour la projection des individus et
			des modalités.}]
		
		\paragraph{Projection des modalités sur le premier plan factoriel}
		
		La projection des modalités sur les deux premiers axes factoriels
		permet de visualiser les associations entre les différentes modalités
		des variables et d'identifier les oppositions structurantes du nuage
		de points.
		
		\begin{figure}[H]
			\centering
			\includegraphics[width=\textwidth]{./img/acm_projection_modalites.png}
			\caption{Projection des modalités sur le premier plan factoriel
				de l'ACM (axes 1 et 2)}
			\label{fig:acm_projection}
		\end{figure}
		
		[Commenter la figure~\ref{fig:acm_projection} en décrivant les
		oppositions observées sur chaque axe. Identifier les modalités
		qui contribuent le plus à la définition des axes et interpréter
		leur signification économique. Exemple~: \textit{L'axe 1 oppose
			les entreprises du secteur tertiaire ayant des besoins en
			financement et en services aux entreprises du secteur primaire
			caractérisées par des besoins en intrants et en équipements.}]
		
		\paragraph{Variables contribuant aux axes factoriels}
		
		Les figures suivantes présentent les contributions des variables
		au premier et au second axe factoriel, permettant d'identifier
		les attributs les plus structurants de chaque dimension.
		
		\begin{figure}[H]
			\centering
			\begin{minipage}{0.48\textwidth}
				\centering
				\includegraphics[width=\textwidth]{./img/acm_contrib_axe1.png}
				\caption{Contributions des modalités au premier axe factoriel}
				\label{fig:acm_contrib1}
			\end{minipage}
			\hfill
			\begin{minipage}{0.48\textwidth}
				\centering
				\includegraphics[width=\textwidth]{./img/acm_contrib_axe2.png}
				\caption{Contributions des modalités au second axe factoriel}
				\label{fig:acm_contrib2}
			\end{minipage}
		\end{figure}
		
		[Commenter conjointement les figures~\ref{fig:acm_contrib1}
		et~\ref{fig:acm_contrib2}. Identifier pour chaque axe les trois ou
		quatre modalités présentant les contributions les plus élevées, et
		interpréter ce que ces contributions révèlent sur la structure
		économique sous-jacente. Les contributions dépassant le seuil moyen
		$\frac{100}{m}$ (où $m$ est le nombre total de modalités) sont
		considérées comme significatives \citep{lebart_morineau_piron_1995}.]
		
		\subsubsection{Classification Ascendante Hiérarchique (CAH)}
		
		La CAH a été conduite à partir des coordonnées factorielles des
		entreprises sur les axes retenus lors de l'ACM, en utilisant la
		méthode d'agrégation de Ward qui minimise la perte d'inertie
		inter-classes à chaque étape de fusion \citep{saporta_2010}. Cette
		approche, combinant ACM et CAH, est particulièrement adaptée aux
		données qualitatives dans la mesure où elle s'affranchit des
		problèmes liés aux distances euclidiennes calculées directement
		sur des variables catégorielles brutes.
		
		\paragraph{Dendrogramme et choix du nombre de classes}
		
		Le dendrogramme obtenu à l'issue de la CAH permet de visualiser
		la hiérarchie des regroupements et d'identifier le nombre optimal
		de classes à retenir.
		
		\begin{figure}[H]
			\centering
			\includegraphics[width=0.85\textwidth]{./img/cah_dendrogramme.png}
			\caption{Dendrogramme de la Classification Ascendante Hiérarchique
				(méthode de Ward)}
			\label{fig:dendrogramme}
		\end{figure}
		
		L'examen du dendrogramme (figure~\ref{fig:dendrogramme}) révèle
		l'existence d'un saut d'inertie particulièrement marqué à l'étape
		correspondant au passage de trois à deux classes, ce qui plaide en
		faveur d'une partition en \textbf{trois classes}. Ce choix est
		conforté par l'éboulis des indices de niveau, qui présente un
		décrochement net à ce niveau. La partition en trois classes offre
		par ailleurs une granularité suffisante pour identifier des profils
		d'entreprises distincts et opérationnellement utiles dans le cadre
		de l'algorithme de mise en relation de la plateforme.
		
		\paragraph{Visualisation de la partition obtenue}
		
		\begin{figure}[H]
			\centering
			\includegraphics[width=\textwidth]{./img/cah_classification.png}
			\caption{Projection des trois classes obtenues par CAH sur le
				premier plan factoriel de l'ACM}
			\label{fig:cah_classes}
		\end{figure}
		
		La figure~\ref{fig:cah_classes} représente les trois classes
		d'entreprises obtenues, projetées sur le premier plan factoriel
		de l'ACM. La séparation visuelle des classes dans cet espace
		factoriel confirme la cohérence de la partition obtenue~: les
		trois groupes occupent des zones distinctes du plan factoriel,
		témoignant d'une bonne homogénéité intra-classe et d'une
		séparation satisfaisante inter-classes.
		
		[Décrire ici les trois classes en termes de profils d'entreprises.
		Exemple de structure~:
		\begin{itemize}
			\item \textbf{Classe 1} --- [Nombre] entreprises (~[X]\,\%)~:
			profil dominant [à compléter à partir des résultats].
			\item \textbf{Classe 2} --- [Nombre] entreprises (~[X]\,\%)~:
			profil dominant [à compléter].
			\item \textbf{Classe 3} --- [Nombre] entreprises (~[X]\,\%)~:
			profil dominant [à compléter].
		\end{itemize}]
		
		\subsubsection{Analyse Factorielle Discriminante (AFD)}
		
		L'AFD est conduite en aval de la CAH afin de valider la qualité de
		la partition obtenue et d'identifier les variables qui discriminent
		le mieux les trois classes d'entreprises \citep{saporta_2010}. Elle
		permet de répondre à deux questions complémentaires~: les classes
		identifiées sont-elles suffisamment distinctes pour être considérées
		comme des groupes homogènes~? Quelles variables permettent de les
		distinguer le plus efficacement~?
		
		\paragraph{Axes discriminants et pouvoir discriminant}
		
		\begin{figure}[H]
			\centering
			\includegraphics[width=0.85\textwidth]{./img/afd_projection.png}
			\caption{Projection des entreprises et des centres de classes
				sur les axes discriminants de l'AFD}
			\label{fig:afd_projection}
		\end{figure}
		
		[Commenter la figure~\ref{fig:afd_projection}. Indiquer le
		pourcentage de variance inter-classes expliqué par le premier
		axe discriminant et par le second. Une valeur cumulée supérieure
		à 80\,\% indique une bonne discrimination. Interpréter la
		séparation des barycentres de classes sur le premier plan
		discriminant.]
		
		\paragraph{Variables discriminantes}
		
		\begin{figure}[H]
			\centering
			\includegraphics[width=0.85\textwidth]{./img/afd_variables.png}
			\caption{Contribution des variables aux axes discriminants de l'AFD}
			\label{fig:afd_variables}
		\end{figure}
		
		[Commenter la figure~\ref{fig:afd_variables} en identifiant les
		deux ou trois variables présentant les contributions les plus
		élevées aux axes discriminants. Ces variables constituent les
		critères les plus pertinents pour la mise en relation des
		entreprises dans la plateforme.]
		
		\paragraph{Taux de bon classement}
		
		La qualité globale de la discrimination est évaluée par le taux
		de bon classement (TBC), calculé par validation croisée afin
		d'éviter le biais de sur-ajustement~:
		
		\begin{equation}
			\text{TBC} = \frac{\text{Nombre d'individus correctement classés}}{n}
			\times 100
			\label{eq:tbc_resultats}
		\end{equation}
		
		[Indiquer ici la valeur du TBC obtenu. Un TBC supérieur à 70\,\%
		est généralement considéré comme satisfaisant dans la littérature
		\citep{lebart_morineau_piron_1995}. Conclure sur la qualité de la
		partition et son adéquation pour alimenter l'algorithme de mise
		en relation de la plateforme.]
		
		\begin{remarque}
			La chaîne analytique ACM $\rightarrow$ CAH $\rightarrow$ AFD
			produit trois résultats directement exploitables par la plateforme
			\textit{Business Partnership}~: \textbf{(i)} les coordonnées
			factorielles de chaque entreprise dans l'espace de l'ACM, qui
			constituent les vecteurs sur lesquels est calculée la similarité
			cosinus pour la mise en relation~; \textbf{(ii)} l'appartenance
			de chaque entreprise à l'une des trois classes, qui permet des
			recommandations de premier niveau~; \textbf{(iii)} les variables
			discriminantes identifiées par l'AFD, qui constituent les critères
			prioritaires de filtrage dans le moteur de recommandation.
		\end{remarque}